% English counterpart of draft/background_and_motivation.md.
% Compile from paper/; see motivation-preview.tex for EuroSys 2027 settings.
\section{Background and Motivation}
\label{sec:background-motivation}

Multi-agent applications complete a user task through interdependent LLM
requests. A worker failure loses both serving capacity and key--value (KV)
state, delaying interrupted requests and potentially their successors.
Backing up KV reduces reconstruction work, but uses memory and bandwidth
that also serve healthy requests. We therefore ask how to allocate limited
protection to reduce workflow completion time (WCT), and when faster
GPU-resident protection earns its additional resource cost.

% One full-width figure with three subfigures; requires subcaption.
% Each landscape panel is 2.24 x 1.85 in at the EuroSys insertion size.
% Plot labels are 10.25 pt at this insertion size; do not shrink this row.
\begin{figure*}[t]
  \centering
  \captionsetup[subfigure]{font=normalsize,labelfont=normalfont,
    textfont=normalfont,justification=centering,singlelinecheck=false}
  \begin{subfigure}[t]{0.32\textwidth}
    \vspace{0pt}
    \includegraphics[width=\linewidth]{figures/motivation-row/motivation-branch-gap.pdf}
    \caption{Workflow-maximum gaps}
    \label{fig:motivation-branch-gap}
  \end{subfigure}\hfill
  \begin{subfigure}[t]{0.32\textwidth}
    \vspace{0pt}
    \includegraphics[width=\linewidth,height=1.85in,keepaspectratio]{figures/motivation-row/motivation-workflow.drawio.pdf}
    \caption{Four-branch workflow}
    \label{fig:motivation-workflow}
  \end{subfigure}\hfill
  \begin{subfigure}[t]{0.32\textwidth}
    \vspace{0pt}
    \includegraphics[width=\linewidth]{figures/motivation-row/motivation-recovery-exposure.pdf}
    \caption{Estimated recovery exposure}
    \label{fig:motivation-recovery-exposure}
  \end{subfigure}
  \caption{Workflow timing and illustrative recovery exposure.
    (a) ECDF of maximum branch-duration gaps per workflow
    (121 workflows; temporally inferred fan-in), not a slack distribution.
    (b) Selected four-branch trace example, with a 91.2\,s gap;
    it is not a maximum-gap sample from (a). Bar positions and lengths
    encode observed timing from the first branch launch; dashed tails
    assume all branches are required at A2's 151.6\,s finish.
    (c) Independent branch-failure counterfactuals at 29.5\,s:
    hatching denotes absorbed delay, solid fill WCT loss.
    Rebuild costs are Qwen3-14B/A100 BF16 static estimates with an assumed
    64-token tail (batch/TP one), not measured failures.}
  \label{fig:motivation}
  \Description{The panels show workflow-maximum duration gaps, a time-scaled
    four-branch fan-in example with a 91.2-second gap, and estimated
    rebuilding cost decomposed into slack-absorbed time and WCT loss.
    A2 determines the assumed join and exposes its full 32.2-second
    rebuilding cost. A1 absorbs its entire cost; A3 and A4 absorb part
    of theirs, leaving WCT penalties of 17.6 and 17.0 seconds.}
\end{figure*}


\subsection{Workflows, Failures, and KV Recovery}
\label{sec:motivation-background}

\paragraph{Multi-agent workflows.}
We represent a workflow as a directed acyclic graph $G=(V,E)$, whose nodes
are LLM requests and whose edges are data dependencies. An agent branch may
include multiple requests and tool calls; a successor runs only after its
dependencies are satisfied. WCT spans the release of the first request to
the completion of the last. For a given execution schedule, a request's
\emph{slack} is the maximum delay to its completion that does not postpone
workflow completion. Requests on a critical path have zero slack.

\paragraph{Serving and failure model.}
We consider workers with block-managed KV, where a worker is a complete
model replica, possibly sharded across a tensor-parallel (TP) group.
A fail-stop GPU or worker failure makes that worker unavailable; its
requests move to a surviving worker with compatible weights already
loaded. The control plane, token history, backup DRAM, and takeover worker
must survive. We address KV recovery under this contract, not whole-host
failures, network partitions, silent corruption, or restoration of lost
model-serving capacity.

\paragraph{Recovery paths and resource costs.}
Without a usable copy, the system must rerun prefill to rebuild lost KV.
A DRAM copy avoids reconstructing the saved prefix but must be transferred
to a GPU before use; concurrent recoveries share transfer bandwidth.
A copy already resident in the actual takeover GPU's HBM can be used
directly, at the cost of HBM capacity that could serve foreground requests
and replication traffic before failure. Backup placement therefore trades
capacity and normal-service overhead against recovery speed.

\paragraph{Relation to existing systems.}
D\'{e}j\`{a}Vu streams replicated KV to remote CPU memory to recover
pipeline execution, spending host storage and transfer
bandwidth~\cite{strati2024dejavu}. GhostServe stores erasure-coded parity
in host memory, trading encoding and reconstruction work for fewer
checkpoint bytes; it also overlaps partial recomputation with recovery
I/O~\cite{jayakody2026ghostserve}. Its evaluated memory-loss recovery
retains a static parallel topology, so recovery of hard device failures
also requires restoring or rebuilding that topology.
LUMEN already coordinates host-memory checkpoint placement and retry
routing by load, and uses draft-model assistance during model reload;
it incurs checkpoint traffic, reserved host memory, and temporary draft
state~\cite{cao2026lumen}. Mooncake pools DRAM/SSD KV storage and schedules
reuse for throughput subject to request-latency SLOs~\cite{qin2025mooncake}.
AegisKV studies a complementary allocation dimension: workflow delay
tolerance together with the KV needed during recovery. Our proposed
distinction is which requests and KV blocks receive scarce fast protection;
neither KV reuse, load-aware recovery, nor transfer overlap is itself new.

\subsection{Recovery Cost Is Not Workflow Delay}
\label{sec:motivation-workflow}

A join waits for its last required branch. Earlier branches can tolerate
some delay, whereas delaying the branch that determines the join time
postpones downstream work. To characterize branch timing differences, we
analyze the public CC/Weka coding-agent
traces~\cite{semianalysis2026cctraces}. From 393 sessions, we first extract
strictly overlapping subagent groups: the latest branch start precedes
the earliest branch finish. We retain groups whose branches all complete
and are followed by a Main request after the last branch finishes.
Main may also issue requests while branches execute. Because the trace
does not record result dependencies, this criterion infers fan-in from
timing; it does not verify consumption of every branch's output.

The filter retains 286 of 292 overlapping groups across 121 workflows.
For each retained group, we compute the longest minus the shortest branch
duration; each workflow contributes its maximum such gap.
Figure~\ref{fig:motivation}(\subref{fig:motivation-branch-gap}) shows a median of
138\,s and a 90th percentile of 1,173\,s. These gaps establish substantial
branch-duration heterogeneity among workflow maxima, not the slack
distribution across requests or failure instants.

Figure~\ref{fig:motivation}(\subref{fig:motivation-workflow}) shows a
selected four-branch example with a 91.2\,s duration gap. Bars preserve
observed starts and durations, aggregating requests and tool calls.
Among eligible groups with four or five branches and a gap within
15\,s of 100\,s, we select the one with the longest critical-branch prefix
to make recovery exposure visible. It is not its workflow's maximum-gap
group. A2 finishes last at 151.6\,s after the fork, and Main issues a
request at 157.4\,s. Under the explicit assumption that Main depends on
all four branches, each branch's margin to A2's completion is its join
slack; the trace supplies times, not these dependency edges.

Figure~\ref{fig:motivation}(\subref{fig:motivation-recovery-exposure}) shows
four independent single-branch failure counterfactuals at 29.5\,s after
the fork, when all four branches have active LLM requests. Trace
hash-block counts give prefixes of 100,864--113,664 tokens. We add an
assumed known 64-token tail and estimate reconstruction time $T_i$ using
AIConfigurator with Qwen3-14B, A100-SXM, BF16, batch one, and TP one.
With join slack $\sigma_i$ derived offline from the observed completions,
and all other execution and queueing unchanged, the WCT penalty is
$\max(0,T_i-\sigma_i)$. This combines trace timing, an assumed dependency,
and substituted static model costs; it is not a failure replay on the
original serving model or a population estimate.

A1--A4 cost 35.2, 32.2, 36.7, and 38.9\,s to rebuild, but expose
0, 32.2, 17.6, and 17.0\,s to WCT, respectively. A1's 91.2\,s slack
absorbs its entire cost. A2 has the lowest rebuilding cost yet the
largest WCT penalty because its slack is zero. Thus, avoided computation
alone can rank protection incorrectly. The example motivates a WCT
objective; it does not establish the accuracy of online slack estimates.

\subsection{Long Contexts and Recovery Tiers}
\label{sec:motivation-long-context}

% Independent single-column Figure 2; native size 3.34 x 2.35 inches.
% CPU-only static estimates; see reuse-aiconfigurator-max-context-h100/.
\begin{figure}[t]
  \centering
  \includegraphics[width=\linewidth]{figures/motivation-row/motivation-recompute.pdf}
  \caption{Static recovery estimates, compared within each model.
    Color/line style denotes path; marker shape denotes model.
    Qwen14B, Llama70B, Qwen235B (2507), and DeepSeek-V3.2 use
    1/4/16/32 H100 GPUs and end at 128K (YaRN)/128K/1,010,000/160K.
    BF16, batch one; total context includes the same 64-token tail.
    HBM holds full history; DRAM assumes 64\,GB/s per GPU without
    overlap; Recompute uses 8K chunks. Hollow markers flag Qwen's
    full-attention projection beyond 256K and DSA table extrapolation
    beyond 32K. K~=~1,024; M~=~$10^6$.}
  \label{fig:motivation-recompute}
  \Description{One plot compares twelve estimated recovery curves on
    logarithmic context and latency axes. Blue solid, orange dashed,
    and red dotted lines encode HBM reuse, DRAM reuse, and Recompute.
    Circles, squares, triangles,
    and diamonds encode Qwen3-14B, Llama-3.1-70B, Qwen3-235B-A22B-Instruct-2507,
    and DeepSeek-V3.2, respectively. The latter two are MoE models;
    DeepSeek additionally uses sparse attention. Curves stop at each
    model's configured context limit. Hollow markers identify Qwen
    full-attention projection beyond 256K and DSA extrapolation beyond
    32K; all values including filled markers are static estimates.
    Full-history HBM is a reference rather than a finite backup policy.
    At 1,010,000 total tokens,
    Qwen 2507 costs an estimated 70 milliseconds for HBM reuse,
    830 milliseconds for DRAM reuse, and 220 seconds for chunked
    full-attention recomputation. GPU allocations differ across models.}
\end{figure}


Repeated agent interactions accumulate code, tool outputs, and other
agents' intermediate results. Backups avoid repeating that computation,
but their location still affects recovery latency.
Figure~\ref{fig:motivation-recompute} compares three initial states:
all historical KV and required auxiliary state reside in the takeover
GPU group's HBM, the same state is in host-local DRAM, or no copy is
available. All paths end after the same known 64-token tail's forward
pass, with history and tail state ready for decoding; weights are
already loaded. HBM reuse therefore still incurs tail computation.

We use AIConfigurator~\cite{nvidia2026aiconfigurator} 0.9.0 HYBRID estimates
for dense Qwen3-14B and Llama-3.1-70B, the Instruct-2507 version of
Qwen3-235B-A22B (235B total, 22B active parameters) with MoE, and DeepSeek-V3.2
(671B total, 37B active) with MoE and DeepSeek Sparse Attention (DSA).
We model Qwen with dense grouped-query attention (GQA) and DSA with its
historical-token indexer. We use
H100-SXM 80\,GiB, vLLM 0.19.0, BF16 weights and KV, and batch one.
TP is 1, 4, 16, and 32, respectively; the MoE models use EP 16 and 32,
and all use PP one. GPU allocations differ, so we compare recovery
paths within each model configuration.

The horizontal axis is total context $C=M+64$, including the known tail.
Each curve reaches its model-specific limit: 131,072 tokens with YaRN
for Qwen3-14B~\cite{qwen2025technical}, 131,072 for
Llama-3.1-70B~\cite{meta2024llama31card}, 163,840 in the official
DeepSeek-V3.2 configuration~\cite{deepseek2025v32config}, and 1,010,000
in Qwen 2507's extended configuration~\cite{qwen2025instruct2507}.
Qwen 2507 natively supports 256K; beyond that length, we project
full-attention costs without Dual Chunk Attention (DCA) position handling
or sparse-kernel acceleration. DSA module tables end at 32K; longer
shapes extrapolate those tables, and cached-prefix queries use analytical
prefix correction. Hollow markers identify these two extrapolation
regimes; filled markers are also static estimates.

HBM reuse computes only the 64-token tail. Recompute starts empty and
processes sequential chunks of at most 8,192 tokens, reusing KV rebuilt
by earlier chunks; vocabulary projection is excluded
from all paths because it does not rebuild KV. Chunking bounds modeled
activation memory; inter-chunk scheduling overhead is excluded.
DRAM first restores each rank's full historical
state, then executes the same tail. GQA transfer bytes account for KV
sharding and head replication. DSA includes replicated MLA latents,
RoPE keys, and native FP8 indexer keys with FP32 scales within the
otherwise BF16 configuration. We assume concurrent restoration at
64\,GB/s per rank, with no shared host bottleneck or transfer--compute
overlap. This bandwidth is a system-model assumption, not a measurement.

At 128K total context, HBM reuse costs an estimated 22--56\,ms,
DRAM reuse 155--357\,ms, and chunked Recompute 14.5--21.0\,s.
At 1,010,000 tokens, Qwen 2507's full-attention cost projection gives
70\,ms, 830\,ms, and 220\,s, respectively. For Qwen3-14B at 128K,
a DRAM copy avoids most historical recomputation; full HBM residency
additionally removes about 335\,ms of modeled transfer waiting.
Thus, backups avoid computation, while HBM's incremental value over
DRAM comes from reducing transfer waiting. Parameter counts, state
representation, and parallelism jointly affect the cost.

All 25 configurations pass the static memory check, with a maximum of
77.8\,GiB per GPU. This checks modeled capacity, not physical deployment,
long-context accuracy, or native sparse recovery correctness.
Full-history HBM residency is an ideal reuse reference: AegisKV's finite
hot-block budget generally leaves some required state in DRAM.
The figure isolates within-model path costs and excludes detection,
queueing, software takeover, and bandwidth contention.

\subsection{Problem and Design Goals}
\label{sec:motivation-goals}

AegisKV asks: \emph{with limited backup capacity and replication and
recovery bandwidth, which KV from which requests should receive what
protection to reduce the workflow-level impact of failures?}
Figure~\ref{fig:motivation-recompute} supports preserving KV, but also
shows why DRAM-only protection is an essential baseline: it avoids most
recomputation without reserving backup HBM. HBM can improve WCT only when
its incremental reduction in recovery waiting is exposed to a workflow.
For an isolated request with slack $\sigma$ and recovery delays
$D\geq H$ under DRAM and HBM reuse, respectively, the saved WCT is
$\max(0,D-\sigma)-\max(0,H-\sigma)$, at most $D-H$ and zero if
$\sigma\geq D$. A large branch-duration gap therefore cannot by itself
justify HBM occupancy, nor can seconds of avoided prefill justify a
subsecond gain over DRAM.

The useful regime must instead combine requests with little remaining
slack, transfer waiting that cannot be hidden by required computation,
and enough spare HBM to avoid harming foreground serving. Concurrent
restores may increase that waiting only when they share a constrained
host or link; that effect is absent from Figure~\ref{fig:motivation-recompute}.
Establishing break-even requires matched DRAM-only and finite-HBM runs
over load, shared bandwidth, backup budget, and failure frequency, counting
both affected and unaffected workflows. Frequency-weighted recovery gains
must justify healthy-service costs from replication and displaced KV.
When they do not, the system should retain DRAM protection and relinquish
the extra HBM. These are validation conditions, not benefits established
by the current static curves.

This problem imposes five goals and constraints.
\textbf{(R1)} Optimize recovery for WCT, accounting for dependency-induced
delay tolerance.
\textbf{(R2)} Allocate limited backup capacity across requests and their
KV to reduce exposed recovery waiting.
\textbf{(R3)} Support concurrent recovery and transfer--compute overlap
without introducing request-level serialization for protection priority.
\textbf{(R4)} Bound normal-service replication overhead and subordinate
backup capacity and bandwidth use to foreground needs.
\textbf{(R5)} Preserve model computation and KV recovery correctness using
existing serving resources.
The next section develops the corresponding placement and recovery design.
